\documentclass[../main.tex]{subfiles}
\begin{document}
%==========================================TIÊU ĐỀ TẬP========================================
\begin{table}[h]
    \begin{adjustwidth}{-1cm}{}
        \begin{tabular}{>{\centering\arraybackslash}p{6cm}|>{\raggedright\arraybackslash}p{12.5cm}}
            \multirow{5}{6cm}
            % Cột bên trái: ÉPISODE và số tập
            {\\[-40pt]
                \flaregothic\fontsize{20pt}{20pt}\selectfont \centering
                \color{eptype}SPÉCIAL\color{black}\\
                \burbank\fontsize{72pt}{72pt}\selectfont
                \color{epnum}XXIII\color{black}\\[6pt]
                \cafeteria\fontsize{20pt}{20pt}\selectfont\color{epnum}(S-23)\color{black}
                \\[18pt]
                % \flaregothic\fontsize{14pt}{14pt}\selectfont
                % \color{epattr}BIENVENUE, \uline{\color{Banpire} Mel} !
                \sen\fontsize{18pt}{18pt}\selectfont
                \color{epattr}Mini-histoire
                \color{black}
            }
             &
            % Dòng thứ nhất: tên tập tiếng Nhật
            {\fotskip\fontsize{18pt}{18pt}\selectfont \ruby{宝}{たから}\ruby{鐘}{かね}マリン。
            }                             \\[6pt]
            % Dòng thứ hai: tên tập tiếng Anh
             & {\cafeteria\fontsize{18pt}{18pt}\selectfont The Stories of a Girl Who Likes Cosplaying Pirates} \\[4pt]
            % Dòng thứ ba: tên tập tiếng Việt
             & {\vnmsans\fontsize{18pt}{18pt}\selectfont Những câu chuyện nhỏ về Houshou Marine}               \\[8pt]
            % Dòng thứ tư: Nhân vật xuất hiện
             & {\caviar{\fontsize{14pt}{14pt}\selectfont Nhân vật: Theo từng câu chuyện}}
            \\[24pt]
        \end{tabular}
    \end{adjustwidth}
\end{table}
%=========================================TRUYỆN CHÍNH========================================
\color{black}
\begin{center}
    \uline{(Đây là những câu chuyện dựa trên các video ngắn (không phải YouTube Shorts) từ kênh của \href{https://www.youtube.com/@HoushouMarine}{Houshou Marine}.)}
\end{center}

\subsection*{\phantomsection\label{ep-S16.I}}
\addcontentsline{toc}{subsection}{{\sen{-Histoire 1-}} {\caviar Con rùa}}
\stpt{-Histoire 1-}
\stcap{Con rùa}
\vspace{-4pt}
{\caviar\fontsize{12pt}{12pt}\selectfont Nhân vật: \emp{kuon1} \emp{flare1} \emp{noel} \emp{marine}}

Bờ biển trải dài dưới ánh hoàng hôn rực rỡ tựa như một bức tranh màu cam vàng được vẽ bằng những tia nắng cuối ngày, phản chiếu lên mặt nước lấp lánh như hàng ngàn viên kim cương. Từng đợt sóng nhẹ nhàng vỗ vào bờ cát, tạo nên một khúc nhạc du dương của thiên nhiên. Gió biển thổi nhẹ, mang theo hơi mặn và một cảm giác bình yên hiếm hoi. Bộ tứ Kuon, Marine, Noel và Flare đang dạo bước, tận hưởng không khí trong lành của buổi chiều muộn.

Thế nhưng, khung cảnh yên bình ấy nhanh chóng bị phá vỡ bởi hành động khó hiểu của nàng Thuyền trưởng.

Marine đột nhiên dừng bước, rồi bò lổm ngổm trên cát như một con vật hoang dã vừa tìm thấy bãi đẻ. Ánh mắt cô long lanh đầy phấn khích, như thể vừa phát hiện ra một bí mật của vũ trụ.

``Em là một chú rùa đây!'' cô hớn hở tuyên bố, rồi ngồi bệt xuống, cố tìm tư thế thoải mái nhất trên mặt cát ấm áp. Cô vung vẩy tay chân trong không khí, trông chẳng khác gì một con rùa biển bị lật ngửa đang cố gắng quay lại.

\img{s-16a2}

Cậu Tổng nheo mắt nhìn cô đầy hoài nghi, tay đưa lên che phần nào ánh nắng chói chang còn sót lại: ``\dots Thật luôn? Em đang làm cái trò gì thế, Thuyền trưởng? Em có biết mình đang ở trên cát chứ không phải ở sở thú không?''

Marine đặt tay lên bụng, nhìn xa xăm, ngó lơ hoàn toàn câu hỏi của cậu. Giọng cô mơ màng như một thi sĩ đang ngẫm nghĩ về kiếp nhân sinh: ``Không biết liệu mình có thể đẻ trứng ở đây không ta\~{}'' cô nói, hai tay vỗ nhẹ lên cát bên cạnh như thể đang chuẩn bị một tổ ấm, ``Trứng rùa biển thì đẻ trên cát mà, đúng không?''

Kuon còn chưa kịp phản ứng thì từ phía xa, một giọng nói khác vang lên, đầy hứng thú như thể vừa nghe thấy một cuộc phiêu lưu mới: ``Ồ? Đang định đẻ trứng hả?''

Cậu quay đầu lại, chỉ kịp thấy Noel đang cầm cây chùy bước tới, nở một nụ cười đầy hứng thú. Đôi mắt cô sáng rực như vừa tìm thấy một trò chơi mới.

``Vậy để Danchou ăn nó nha!''

Nói rồi, nàng Hiệp sĩ bắt đầu đá nhẹ vào lưng Marine, như thể đang thúc giục cô nàng đẻ trứng thật sự. Mỗi cú đá đều đi kèm với một tiếng hô vui vẻ: ``Hây! Hây, hây hây! Nào, đẻ trứng nhanh lên nào! Rùa mà đẻ chậm là bị sóng cuốn đi đấy!''

Senchou lăn lộn trên cát, cố bò ra xa, nhưng Noel vẫn cứ đuổi theo với sự nhiệt tình của một người đang tham gia một trò chơi thể thao: ``\textbf{Á á á!! Từ từ đã!! Tôi không phải rùa thật đâu!!}'' Cô nàng hét lên, nhưng giọng cô vẫn đầy vẻ phấn khích, như thể cô đang tận hưởng từng giây phút của trò chơi này.

Kuon chỉ biết đứng đó, mặt đơ ra không hiểu chuyện đang xảy ra trước mặt. Cậu đưa tay xoa trán, tự hỏi tại sao mình lại nghĩ rằng một buổi dạo biển sẽ yên bình: ``\dots Mấy đứa làm cái gì vậy? Em ấy nói là rùa, thế em liền đá em ấy như một quả bóng?''

Noel vẫn vô cùng hăng hái, vừa đá vừa giải thích với giọng của một chuyên gia động vật học: ``Thúc giục rùa đẻ trứng đó mà, Kuon-senpai! Ở ngoài thiên nhiên, người ta phải giúp rùa bằng cách kích thích chúng! Em đọc trên Internet đấy!''

Trong lúc nàng Hiệp sĩ đang hành hạ nàng Thuyền trưởng một cách khó hiểu, Flare bước đến từ phía sau, khoanh tay nhìn cảnh tượng trước mặt. Cô lắc đầu, nhưng khóe miệng lại giấu không được một nụ cười: ``Cậu ấy lại bị bắt nạt nữa rồi\dots\ Mà trông cậu ấy có vẻ thích thế.''

\img{s-16a1}

Cậu Tổng liếc nhìn Marine, người đang cười khúc khích giả tiếng khóc trong khi vừa bò vừa né mấy cú đá nhẹ từ Noel, rồi lại quay sang nhìn nàng Yêu tinh vàng. Cậu nhún vai: ``\dots  Anh thấy hình như em ấy muốn như vậy đấy. Em ấy máu khổ dâm mà. Nếu không thích, em ấy đã chạy từ lâu rồi, chứ không nằm đó mà rên rỉ như một con rùa bị bắt nạt.''

Flare bật cười, kéo tay Kuon ra xa khỏi ``hiện trường'' hỗn loạn: ``Ra đây để em kể cho anh nghe câu chuyện đằng sau vụ này. Hồi đó bọn em đang livestream về động vật biển, và Marine đã đùa là muốn làm rùa. Từ đó Noel luôn coi đó là nhiệm vụ của mình.''

Và thế là hai người họ cứ thế bỏ mặc nàng Thuyền trưởng, để cô tiếp tục ``chịu trận'' dưới tay của Noel. Họ đi dọc bờ biển, thỉnh thoảng lại nghe thấy tiếng cười và tiếng la hét vọng lại.

Dưới ánh chiều tà, một ``con rùa'' tóc đỏ vẫn đang lăn lóc trong ``đau đớn'' (?) trên bãi cát, trong khi một nàng Hiệp sĩ vẫn tiếp tục đá cô như một quả bóng. Tiếng cười của họ hòa lẫn với tiếng sóng, tạo nên một bản nhạc đặc trưng của \hll - vừa hỗn loạn, vừa ấm áp.

Cậu nhìn về phía hai người họ, rồi quay sang Flare: ``Anh chỉ hy vọng là em ấy không bắt đầu đẻ trứng thật. Công ty không có ngân sách cho một con rùa thứ hai.''

Flare cười, vỗ vai tôi: ``Yên tâm, em ấy chỉ đóng kịch thôi. Mà nếu có đẻ, chắc cũng chỉ là trứng sô-cô-la mà Marine giấu trong túi.''

Cậu thở dài, vừa cười vừa lắc đầu: ``\dots Chả hiểu nổi em ấy đang nghĩ gì nữa\dots''

\pagebreak

\subsection*{\phantomsection\label{ep-S16.II}}
\addcontentsline{toc}{subsection}{{\sen{-Histoire 2-}} {\caviar Sô-cô-la ngày lễ Tình nhân}}
\stpt{-Histoire 2-}
\stcap{Sô-cô-la ngày lễ Tình nhân}
\vspace{-4pt}
{\caviar\fontsize{12pt}{12pt}\selectfont Nhân vật: \emp{kuon1} \emp{fubuki} \emp{mio} \emp{okayu} \emp{korone} \emp{marine}}

\pov{Hikawa Kuon}

Một ngày Lễ tình nhân nọ, không khí trong văn phòng ngập tràn sắc hồng và những trái tim giấy treo lủng lẳng. Tôi đang ngồi làm việc riêng cùng Korone - cô bạn chó đang say sưa với một gói kẹo mới mua - thì cả hai bất chợt nghe thấy một giọng nói quen thuộc gọi lớn từ phía hành lang, vang dội như tiếng chuông báo động: ``\textbf{Korone-senpai! Kuon-senpai!}''

Chưa kịp để họ đáp lời, Marine đã xuất hiện trước mặt với một nụ cười rạng rỡ như ánh nắng đầu xuân. Trên tay cô là một hộp sô-cô-la hình trái tim màu đỏ, được gói trong giấy bóng kính lấp lánh, trông như một báu vật vừa được khai quật.

``Mừng Lễ tình nhân vui vẻ!''

Tôi đứng ngây ra một lúc, nhìn hộp quà, rồi nhìn Marine, rồi lại nhìn hộp quà. Một khoảnh khắc lưỡng lự, nhưng cuối cùng tôi cũng gật đầu: ``Ờm\dots\ cảm ơn\dots?'' giọng tôi vẫn còn hơi ngập ngừng, không biết nên phản ứng thế nào với một món quà từ Thuyền trưởng trong ngày Valentine.

Korone thì mắt sáng lên, nhanh chóng nhận lấy hộp quà với vẻ mặt của một đứa trẻ vừa được tặng đồ chơi mới: ``Cảm ơn em nha, Maririn! Em đúng là thiên thần!''

Thấy thế, tôi cũng nhún vai, đáp với giọng tỉnh bơ nhưng không giấu được chút ngạc nhiên: ``Em tặng thì anh xin. Anh cũng vừa mới nhận sô-cô-la từ Korone-chan nữa kìa. Hôm nay hình như là ngày của những món ngọt.''

Sau màn trao quà đầy hứng khởi, họ quyết định ăn luôn tại chỗ, không để món quà nguội đi. Korone vui vẻ hô lên, tay giơ cao hộp sô-cô-la như một chiến lợi phẩm: ``Mời mọi người nghen\~{}!''

Tôi cũng gật đầu, bẻ một miếng bỏ vào miệng. Vị ngọt tan chảy trên đầu lưỡi: ``Ừ ăn nào. Từ sáng đến giờ anh đã ăn gì đâu, chỉ có cà phê với bánh mì khô thôi.''

Tôi nhai ngon lành, vì vốn dĩ sô-cô-la là một trong những món tôi thích - nó giúp tôi quên đi cơn đói và cả cơn đau đầu từ đống báo cáo. Korone cũng hăng hái ăn theo, nhai nhồm nhoàm với vẻ mặt mãn nguyện của một người vừa tìm thấy kho báu. Nhưng chỉ vài giây sau, cô bỗng khựng lại, mắt trợn tròn như hai viên bi, rồi ngã ngửa ra sàn với một tiếng *RẦM* nặng nề. Cục sô-cô-la mà cô cầm trên tay bị bật lên không trung, xoay vài vòng rồi rơi xuống cạnh mặt cô.

``\textbf{Korone-senpai!!}'' Marine hoảng hốt kêu lên khi thấy đàn chị đổ thẳng xuống đất, mặt cô tái xanh như thể vừa nhìn thấy một vụ án, ``Có chuyện gì vậy!? Có cần gọi cấp cứu không!?''

Tôi cũng tá hỏa, đặt hộp sô-cô-la xuống, lao tới bên cạnh Korone: ``Có làm sao không vậy!? Sao tự dưng lại lăn đùng ra ngỏm rồi!? Em ấy có bị dị ứng gì không?''

Trong khi cả hai đang hoang mang, Korone lẩm bẩm một câu tiếng Anh bồi, miệng vẫn còn ngậm cục sô-cô-la, giọng đều đều như một cái máy ghi âm đang phát bản ghi cuối cùng: ``\eng{Ồ\dots\ tôi chớt rồi. Mãi mãi cảm ơn\~{}}\footnote{Nguyên văn: ``Oh, I'm die. Thank you forever.'' Đây là câu nói ``xai nữ phát'' nổi tiếng của Korone, bắt nguồn từ buổi phát sóng trò chơi \textit{KTANE} (\textit{Keep Talking and Nobody Explode}) với Mori Calliope (sẽ được giới thiệu trong quyển số Mười Sáu) vào tháng 11 năm 2020.}''

\img{s-16b2}

Ngay lúc ấy, họ mới nhận ra vấn đề. Cả hai người xanh xao mặt mày, hoảng hốt đồng thanh, như thể vừa phát hiện ra một sự thật hiển nhiên mà lại quên mất: ``À đúng rồi, chó có ăn được sô-cô-la đâu! Chúng có độc mà!''

Marine run lên cầm cập, lắp bắp, tay bấm số điện thoại cấp cứu nhưng ngón tay không ngừng run: ``K-Không thể nào\dots\ Tại sao chứ!? Em tưởng chị ấy sẽ ổn vì chị ấy là người chó mà\dots''

Tôi quay sang nhìn cô, giọng đầy bất lực, tay xoa trán: ``Tự dưng lại cho em ấy ăn cái đó làm gì! Em ấy là chó--''

Bất chợt tôi dừng lại, bất chợt bình tĩnh tới lạ thường. Một suy nghĩ lóe lên trong đầu: ``À mà, có gì đó sai sai\dots''

Marine nhíu mày, vẫn còn đang hốt hoảng tìm điện thoại gọi cấp cứu, giọng đầy biện minh như một nhà khoa học đang bảo vệ giả thuyết của mình: ``Em tưởng chị ấy sẽ ăn được vì chị ấy là người chó mà phần người lại nhiều hơn phần chó chứ!''

\begin{wrapfigure}[18]{l}{0pt}
    \begin{tikzpicture}[scale=4]

        % Bán kính hình tròn
        \def\radius{1}

        % Vẽ các phần của biểu đồ
        % Phần hồng (góc 0 -> 35)
        \draw[fill=pink!70] (0,0) -- (0:\radius) arc (0:35:\radius) -- cycle;
        % Phần nâu (góc 35 -> 80)
        \draw[fill=brown!70] (0,0) -- (35:\radius) arc (35:80:\radius) -- cycle;
        % Phần vàng (góc 80 -> 360)
        \draw[fill=yellow!70] (0,0) -- (80:\radius) arc (80:360:\radius) -- cycle;

        % Viền ngoài hình tròn
        \draw (0,0) circle (\radius);

        % Thêm nhãn cho các phần
        % Góc giữa phần hồng là 0 -> 35, nên midpoint là 17.5
        \node at (17.5:0.6) {Bí ẩn\\(9.7\%)};
        % Góc giữa phần nâu là 35 -> 80, midpoint là 57.5
        \node at (57.5:0.6) {Chó\\(12.5\%)};
        % Góc giữa phần vàng là 80 -> 360, midpoint xấp xỉ 220
        \node at (220:0.6) {Người\\(77.8\%)};
    \end{tikzpicture}
    \caption{Biểu đồ thành phần của Inugami Korone, dựa theo suy luận của Marine.}
\end{wrapfigure}

Tôi ngẫm lại một lúc lâu, và thở dài: ``\dots  Nói cũng có lý. Nhưng mà em có thể giải thích phần ``bí ẩn'' đó là gì không?''

Ngay khi vừa dứt lời, Korone bỗng ngồi bật dậy, mắt long lanh như thể chưa có chuyện gì xảy ra. Cô vẫn nhai cục sô-cô-la, và cười toe: ``À, ừ, đúng rồi ha! Chó không ăn được sô-cô-la, nhưng mà em là người chó, nên em vẫn sống!''

Cả hai người thở phào nhẹ nhõm, nhưng rồi đồng thanh bất bình thốt lên, giọng như một bản hòa ca của sự bực bội: ``\textbf{\dots\ Làm hết hồn bọn này rồi!}''

Bất chợt, cánh cửa phòng mở ra, và các thành viên còn lại của \textit{hololive GAMERS} bước vào. Mio vẫy tay chào, tươi cười hỏi, như thể vừa nhận ra ngày hôm nay đặc biệt: ``Chào buổi sáng! Hôm nay là ngày Lễ tình nhân đấy! Có ai tặng quà cho chị không?''

Fubuki hí hửng nhìn Marine, hai tay cô tạo dáng vuốt mèo, mắt sáng rực: ``Marine-chan còn sô-cô-la nào cho chị không đây? Chị đã sẵn sàng để nhận quà từ em rồi.''

Okayu thì lười biếng nhưng không quên đưa tay ra, giọng đều đều nhưng đầy vẻ háo hức: ``Em cũng mang cho chị sao? Cảm ơn nha, Thuyền trưởng\~{} Em đã chuẩn bị sẵn một cái bụng trống để ăn rồi đấy.''

Nghe thấy tiền bối mình đến để nhận quà, Marine hào hứng đáp, giọng như một người bán hàng đón khách: ``À, tất nhiên em cũng mang đến tặng cho các chị đây! Sô-cô-la Valentine dành cho tất cả mọi người!''

\img{s-16b4}

Tôi chỉ đơn giản nhún vai, nói hai chữ, vẫn ngồi nhấm nháp miếng sô-cô-la của mình: ``Cứ việc.''

Nhưng điều buồn cười là ngay sau khi họ ăn sô-cô-la, cả ba người cũng\dots\ ``thăng thiên'' y hệt Korone ban nãy, mỗi người một biểu cảm đầy hài hước.

Fubuki ngã ngửa ra sau, hai tay duỗi thẳng, mắt trợn tròn như thể vừa thấy Chúa. Mio thì gục xuống bàn, tay vẫn cầm miếng sô-cô-la như một người vừa bị trúng đòn chí mạng. Okayu thì nhẹ nhàng hơn - cô chỉ đơn giản là nằm lăn ra ghế, như một con mèo đã ngủ quên giữa giấc mơ ngọt ngào.

Lý do? Các bạn cứ nhìn bức ảnh ở dưới là biết. Hóa ra, sô-cô-la độc với chó và mèo, và cả ba đều là người thú với phần thú khá đáng kể.

\img{s-16b5}

Tôi nhìn cảnh tượng trước mặt, tay chống cằm, lẩm bẩm với giọng đầy ngạc nhiên: ``Ơ? Anh tưởng mấy đứa cũng là người thú với phần người hơn phần thú như Koro-chan chứ? Hay là YAGOO đã bí mật tăng tỷ lệ thú của các em lên?''

Marine thì chỉ biết gào lên đầy bất lực, giọng xé toạc không khí: ``\textbf{Tại sao GAMERS các chị lại toàn chó với mèo thế kia dzậy!? Tôi đã tặng sô-cô-la cho toàn bộ nhóm GAMERS và tất cả đều ngã lăn ra!}''

Tôi thở dài, nhún vai đáp lại, với giọng của một người đã quá quen với sự vô lý của công ty này: ``Chuyện đó thì\dots\ chỉ có YAGOO mới biết em ạ. Đừng hỏi. Có thể ông ấy đã cố tình tạo ra một nhóm người thú để trêu chúng ta.''

\pagebreak

\subsection*{\phantomsection\label{ep-S16.III}}
\addcontentsline{toc}{subsection}{{\sen{-Histoire 3-}} {\caviar Tình yêu và phản bội}}
\stpt{-Histoire 3-}
\stcap{Tình yêu và phản bội}
\vspace{-4pt}
{\caviar\fontsize{12pt}{12pt}\selectfont Nhân vật: \emp{kuon1} \emp{rushia} \emp{marine}}

Một buổi chiều tưởng chừng bình thường tại văn phòng \hll, tôi đang cặm cụi với đống giấy tờ thì một giọng nói quen thuộc vang lên từ ngoài cửa, đầy vẻ tươi tắn: ``Chào mọi người-nanodesu!''

Rushia lấp ló ngoài cửa, đôi mắt đỏ đảo quanh căn phòng một cách tò mò, như một chú chim non vừa bay đến tổ mới. Cô đang tìm kiếm ai đó để trò chuyện, nhưng ngay khi ánh nhìn lướt qua, cô lập tức bắt gặp một cảnh tượng kỳ lạ: Marine đứng khoanh tay, quay lưng về phía cửa sổ, dáng vẻ như đang suy tư điều gì đó rất nghiêm túc - kiểu suy tư của một thuyền trưởng đang tính toán hải trình, chứ không phải của một người bình thường.

``Marine-senchou?'' Rushia cất tiếng gọi, giọng vừa tò mò vừa có phần dè dặt.

Nghe thấy giọng cô, nàng Thuyền trưởng từ từ xoay người lại. Nhưng thay vì đáp lời như bình thường - kiểu ``Ahoy!'' hay ``Có chuyện gì thế?'' - Marine đột nhiên tạo dáng đầy gợi cảm. Cô lắc nhẹ hông, ưỡn người ra phía trước, quàng hai tay sau gáy như một người mẫu đang tập tạo dáng cho tạp chí thời trang. Đôi mắt cô ánh lên hình trái tim khi cô cất giọng đầy quyến rũ, như một nhân vật trong phim tình cảm Pháp: ``\eng{I'm horny\~{}}''

\img{s-16c1}

\dots Im lặng.

Rushia cứng đờ, mắt mở lớn, miệng há hốc. Trong đầu cô chỉ có một suy nghĩ duy nhất, như một dòng chữ chạy trên bảng quảng cáo: ``\textbf{\textit{Có kẻ kỳ quặc trong phòng kìa!}}'' Cô muốn hét lên, nhưng lại không dám thốt ra thành lời, chỉ có thể đứng yên như một pho tượng.

Cô chớp mắt vài lần, cố gắng hiểu chuyện gì đang diễn ra trước mặt. Một phần trong cô tự hỏi: ``\textit{Không lẽ mình vào nhầm phòng rồi? Hay đây là một buổi quay phim nào đó?}''

Cô khẽ lùi lại một chút, định quay người bỏ đi, nhưng ngay lúc đó, Marine đột nhiên quay người lại, đôi mắt rực sáng như thấy được kho báu dưới đáy biển. Cô nhận ra Rushia đang ở đó, và phản ứng của cô giống như một kẻ nghiện game vừa thấy tựa game mới ra mắt.

``\textbf{Rushia!!!}''

Rushia giật mình, đứng sững tại chỗ. Chưa kịp phản ứng, cô thấy Marine đang tiến lại gần với ánh mắt si mê và một nụ cười tràn đầy ý đồ đáng ngờ - kiểu nụ cười mà bạn chỉ thấy trong các bộ phim kinh dị khi kẻ điên sắp ra tay.

Rushia bắt đầu cảm thấy có gì đó không ổn, lùi dần về phía cửa: ``Đ-đợi đã, hình như có gì đó sai sai\dots\ Senchou, bà có ổn không? Có cần tôi gọi bác sĩ không?''

Cô chưa kịp nói dứt câu, Marine đã lao đến với tốc độ đáng sợ - nhanh hơn cả lúc cô đang chạy trốn khỏi những trò đùa của Noel - rồi dí sát vào người cô, bày tỏ sự thắm thiết như một cô gái đang tỏ tình trong phim lãng mạn Hàn Quốc: ``Rushia\~{}, hun tui, hun tui đi\~{}! Trông bà đáng yêu lắm á! Yêu bà nhứt lun\~{} Tui không thể sống thiếu bà được!''

\img{s-16c2}

Tôi cũng vừa tới văn phòng sau khi đi lấy một đống tài liệu. Thấy cảnh tượng có phần khó đỡ kia, tôi đứng đơ người. Nhìn cảnh tượng Marine lao tới ôm ấp, tiếp cận, gạ gẫm Rushia với thái độ khó hiểu và có phần khó chịu. Tôi vẫn còn cầm tập giấy trên tay, và tôi đã suýt làm rơi nó.

``\dots  Em đang làm cái quái gì thế, Senchou? Em có biết mình vừa nói gì không?''

Rushia lập tức thoát khỏi vòng vây của đồng đội ``bỗng lên cơn'', phồng má tức giận và hét lên, giọng như một người vừa phát hiện ra một bí mật đen tối: ``Im miệng đi, cái đồ ngoại tình này!''

Tôi ngơ người, mặt đầy dấu hỏi. Một cơn gió lạnh thổi qua, và tôi tự hỏi liệu mình có bị lạc vào một vở kịch nào đó không: ``\dots  Hả? Ngoại tình? Với ai?''

Không để tôi kịp hiểu chuyện gì, Marine nhào đến, hét lớn như một vị thuyền trưởng ra lệnh cho thủy thủ đoàn trong một cơn bão, giọng đầy quyền uy: ``\textbf{Câm miệng! Quỳ xuống và kết hôn với ta đi! Đây là mệnh lệnh của Thuyền trưởng đó, có biết không hả!?}''

\img{s-16c3}

Cả hai người đồng loạt giật mình. Rushia mắt mở to, còn tôi thì đứng chết trân, không biết nên cười hay nên khóc.

Tôi liền hạ giọng, cố gắng dỗ dành Marine với giọng của một người đang cố trấn an một đứa trẻ đang lên cơn: ``Nói như dọa thế kia người ta sợ em là phải\dots\ Em có thể từ từ nói không? Bình tĩnh đi, Thuyền trưởng.''

Nhưng thật bất ngờ, thay vì phản ứng giận dữ hay lúng túng, Rushia lại đứng đờ người, mắt chớp chớp mấy lần rồi lẩm bẩm trong vô thức. Đôi mắt cô hiện rõ hình trái tim, như thể cô đang bị một loại bùa chú nào đó chi phối: ``Ể? Đ-Được thôi\dots''

Tôi trố mắt, suýt làm rơi tập giấy trong tay: ``\dots  Cái quần què gì đang xảy ra thế\dots? Ủa, l-là sao? Em đồng ý thật á? Em có biết mình vừa nói gì không?''

Tôi còn chưa kịp hiểu chuyện gì, thì Marine đã nhanh tay lấy điện thoại ra, bấm số gọi ngay cho ai đó với vẻ mặt đầy phấn khích của một người vừa chiến thắng: ``Alo? Fuu-tan đó hả? Em có tin được không, Rushia vừa đồng ý kết hôn với em! Em sẽ cưới cô ấy vào cuối tuần này!''

Nhưng chưa kịp nói thêm, cô đã ăn ngay một cú đấm trời giáng từ Rushia kèm theo tiếng hét xé trời, đủ để làm rung cả tòa nhà: ``\textbf{CÁI THẰNG KIA!!!!!!!!!!!}''

\img{s-16c4}

Tôi đứng đó, chỉ biết thở dài, nhìn cảnh Marine bay ngược ra sau: ``Không biết phải nói thế nào, nhưng anh cảm thấy em cũng đáng bị vậy đấy\dots\ Tự nhiên đòi cưới người ta mà không hỏi ý kiến.''

Chẳng ngờ rằng, cú đấm của Rushia quá mạnh - mạnh hơn cả những cú đấm của cô trong các buổi stream - đến nỗi Marine bị hất văng khỏi văn phòng, phá tung cả cửa sổ. Tôi có thể thấy bóng dáng cô bay qua khung cửa, tay chân vung vẩy trong không trung như một con búp bê vải.

``Aaaaaaa\textbf{aaaaaah!!!}''

*XOẢNG!*

Rushia phủi bụi trên tay một cách thản nhiên, như thể vừa hoàn thành một nhiệm vụ hàng ngày, rồi lấy điện thoại ra làm việc riêng. Khuôn mặt cô bày tỏ sự thỏa mãn, giống như một người vừa giải quyết xong một vấn đề nan giải: ``Đáng đời kẻ bắt cá hai tay! Tưởng tôi dễ tính lắm à?''

Trong thâm tâm của cô nàng Pháp sư, cô thà trả tiền bồi thường cửa kính vỡ còn hơn là cặp kè (dù một phần cô thích vậy) với Thuyền trưởng. Ít nhất thì cửa kính không đòi hỏi cô phải hứa hẹn gì.

Còn tôi? Tôi chỉ đứng nhìn đống đổ nát và chép miệng, lắc đầu khi thấy những mảnh kính vỡ lấp lánh trên sàn: ``Lại chuẩn bị mất tiền mua kính rồi\dots\ Tuần này là lần thứ ba rồi đấy. Ai đó nên nói với Marine là đừng có tán tỉnh người ta theo kiểu đó nữa.''

Tôi thở dài, cầm điện thoại lên gọi cho thợ sửa kính, và thầm nghĩ: ``\textit{Ít nhất thì lần này không có ai bị thương nặng. Ngoại trừ lòng tự trọng của Marine.}''

\pagebreak

\subsection*{\phantomsection\label{ep-S16.IV}}
\addcontentsline{toc}{subsection}{{\sen{-Histoire 4-}} {\caviar Cú tát}}
\stpt{-Histoire 4-}
\stcap{Cú tát}
\vspace{-4pt}
{\caviar\fontsize{12pt}{12pt}\selectfont Nhân vật: \emp{kuon1} \emp{rushia} \emp{marine} \emp{lamy}}

\begin{center}
    \uline{(Bắt nguồn từ meme Jack Sparrow bị tát trong \textit{Cướp biển vùng Caribe}.)}
\end{center}

Trong văn phòng, không khí đang toát lên một vẻ yên tĩnh hiếm có - cái kiểu yên tĩnh mà bạn chỉ có được khi không có Coco hay Kanata ở gần. Tôi ngồi ở bàn làm việc, mắt dán vào màn hình máy tính, ngón tay lướt nhẹ trên bàn phím, tận hưởng một khoảnh khắc bình yên quý giá. Bên ngoài cửa sổ, nắng vàng trải dài, và dường như cả thế giới đang im lặng để tôi hoàn thành nốt báo cáo cuối cùng.

Bỗng nhiên, cánh cửa bật mở một cách mạnh bạo, như thể có ai đó vừa đá tung nó. Cánh cửa va vào tường với một tiếng *RẦM* lớn, và một bóng người xuất hiện.

Lamy bước vào với vẻ mặt cau có, hai tay siết chặt như thể đang cầm một thứ gì đó đầy giận dữ. Cô hậm hực tiến thẳng đến Marine, mà không thèm để ý đến ai khác. Đôi mắt cô như muốn phóng hỏa, và bước chân của cô vang lên những tiếng nặng nề trên sàn.

\img{s-16d1}

Marine ngước lên, chớp mắt một cái rồi nở nụ cười xã giao, tưởng rằng Lamy đến để chào hỏi bình thường: ``A, Lamy--''

*CHÁT!*

Một cái tát giáng thẳng vào má nàng Thuyền trưởng, mạnh đến nỗi cô giật mình lùi một bước. Âm thanh của cú tát vang vọng khắp căn phòng, như một tiếng súng nổ trong bầu không khí yên tĩnh. Tôi, đang theo dõi sự việc từ đầu đến cuối, khẽ nhíu mày, tay ngừng gõ phím.

``Lại bị ăn tát à?'' tôi hỏi, giọng nửa thắc mắc, nửa buồn cười. ``Em vừa làm gì mà bị người ta tát thế?''

Marine đưa tay xoa má, miệng méo xệch, mắt vẫn còn mở to vì sốc: ``Chả biết em có đáng bị như thế không nữa\dots\ Em chẳng làm gì cả!''

Tôi nhún vai, giọng như một người đã chứng kiến quá nhiều chuyện: ``Chắc phải làm gì mới để người ta tát chứ. Không ai tự dưng tát người khác đâu.''

Marine thở dài, vừa định quay đi thì cánh cửa văn phòng lại mở lần thứ hai. Cô quay đầu lại, và thấy Rushia bước vào với vẻ mặt đầy nghi hoặc.

``Ồ, Rushia!'' Marine vui vẻ gọi, nhưng chưa kịp nói gì thêm thì nàng Pháp sư tóc xanh đã sải bước đến gần, mắt nhìn Marine với ánh mắt như thể vừa phát hiện ra một kẻ phản bội.

``Vừa rồi là ai vậy?'' nàng Pháp sư hỏi, giọng điệu đầy nghi hoặc, như đang chất vấn một tội phạm.

Marine chưa kịp phản ứng, chỉ vừa hé môi: ``Hả?''

*CHÁT!*

\img{s-16d2}

Cô lại ăn thêm một cú tát trời giáng từ Rushia còn mạnh hơn cả cú của Lamy, đến nỗi má cô ửng đỏ lên rõ rệt và suýt nữa thì cô ngã nhào. Rushia tát xong cũng chẳng nói thêm câu nào, chỉ quay lưng bước đi như thể chưa từng có chuyện gì xảy ra, để lại Marine đứng đờ ra như một pho tượng.

Tôi chậm rãi gật gù, giọng đầy thương cảm nhưng cũng không giấu được chút hài hước: ``Ái da\dots\ Trông đau đấy. Hai cú tát liên tiếp - kỷ lục mới của em đấy, Thuyền trưởng.''

Marine đứng đờ ra vài giây, tay xoa má, mắt nhìn theo Rushia đã khuất, rồi lẩm bẩm như vừa ngộ ra điều gì đó, giọng đầy vẻ đầu hàng số phận: ``\dots  Cái này thì em đáng bị như thế thật. Em nghĩ em biết tại sao rồi.''

\img{s-16d3}

Tôi khoanh tay, gật gù với vẻ mặt của một người đang đưa ra lời khuyên: ``Có khi em phải ăn thêm hơn chục cái nữa thì mới đủ. Tùy vào tội của em là gì.''

Cô nhìn tôi bằng ánh mắt khó hiểu, nhưng tôi chỉ nhếch môi cười trừ rồi ra hiệu, quay lại với công việc: ``Không có gì đâu. Anh chỉ đùa thôi. Nhưng nếu em muốn chuộc tội, anh nghĩ em nên bắt đầu bằng cách mua trà sữa cho cả phòng.''

Marine chỉ biết lắc đầu, vừa cười vừa thở dài, tay vẫn xoa má. Trên mặt cô, hai vệt đỏ rõ ràng in hằn như một lời nhắc nhở về những gì vừa xảy ra.

\pagebreak

\subsection*{\phantomsection\label{ep-S16.V}}
\addcontentsline{toc}{subsection}{{\sen{-Histoire 5-}} {\caviar Xe ba bánh}}
\stpt{-Histoire 5-}
\stcap{Xe ba bánh}
\vspace{-4pt}
{\caviar\fontsize{12pt}{12pt}\selectfont Nhân vật: \emp{kuon1} \emp{kanata} \emp{marine}}

\pov{-}

Một buổi sáng đầy nắng, những tia nắng vàng óng tràn qua khung cửa sổ, đánh thức nàng Thiên sứ khỏi giấc ngủ. Kanata vừa dụi mắt vừa lết ra khỏi giường, đầu óc vẫn còn mơ màng. Cô nhìn đồng hồ, và ngay lập tức, mắt cô mở to, giọng hét vang khắp căn phòng: ``\textbf{Ôi không! Trễ mất rồi!}''

Kanata vội vã chộp lấy bộ đồng phục, mặc qua loa trong vòng chưa đầy ba mươi giây, rồi lao ra khỏi nhà như một cơn lốc. Trên đường chạy, cô thở dốc, mồ hôi lấm tấm trên trán, trái tim đập thình thịch như thể đang trong một cuộc thi chạy marathon.

``Chết tiệt\dots! Để xem nào\dots'' cô nghĩ bụng, rồi giang đôi cánh trắng muốt ra, định bay thẳng đến văn phòng để tiết kiệm thời gian. Nhưng vừa vỗ cánh một cái - *PHỤT!* - đôi cánh co giật một hồi rồi bỗng nhiên\dots\ quăn lại như một tờ giấy bị ướt, không còn chút sức mạnh nào.

Cô hoảng loạn chạm vào cánh mình, mắt mở to đầy sợ hãi: ``Thôi chết! Cánh của mình bị hư rồi! Làm sao mà bay được nữa!''

\img{s-16e1}

Vậy là khỏi bay. Giờ chỉ còn cách chạy bộ thôi - và với đôi chân ngắn của một thiên sứ, việc chạy đến văn phòng sẽ mất ít nhất ba mươi phút.

``Kiểu này mình sẽ bị muộn mất\dots'' Kanata rên rỉ, cố hết sức lao đi trong khi đồng hồ vẫn không ngừng nhích tới từng giây.

Nhưng đúng lúc đó, cô bỗng thấy một bóng người thấp thoáng ở phía xa. Một chiếc xe đang di chuyển, nhưng không phải là một chiếc ô tô, cũng chẳng phải là xe máy, mà là\dots\ một chiếc xe ba bánh! Nó màu đỏ, trông như một món đồ chơi khổng lồ dành cho trẻ em, nhưng lại đang di chuyển trên đường với tốc độ đáng kinh ngạc.

``Hả\dots?'' Kanata cảm thán, cố nhìn cho rõ hơn.

Và người đang ngồi trên chiếc xe ấy không ai khác chính là Marine, nàng Thuyền trưởng của \textit{hololive}. Nhưng thay vì trên một chiếc xe ô tô hoặc xe nào đó mà người thường hay đi, cô lại\dots\ đang cưỡi một chiếc xe đạp ba bánh trẻ con!

\img{s-16e2}

Nàng ``Hải tặc'' đạp xe một cách đầy tự hào, vừa lái vừa huýt sáo, trông như một đứa trẻ mới sắm được đồ chơi mới. Khi thấy Kanata, cô dừng xe lại, vẫy tay nhiệt tình, giọng đầy hứng khởi: ``Ê, Kanatan! Đang vội hả? Leo lên xe đi, chị chở! Chiếc xe này chạy nhanh lắm đấy!''

Kanata đứng im một lúc, nhìn chằm chằm vào chiếc xe kỳ quặc kia. Nó không phải một chiếc xe bình thường: phần bánh trước thì quá to so với hai bánh sau, còn cái yên xe thì bé tí, trông chẳng đáng tin chút nào. Chưa kể đến việc ngồi trên xe dành cho con nít còn khiến người ngoài nhìn vào chê cười.

Cô im lặng, rồi\dots\ quay mặt đi. Không nói một lời, cô vỗ đôi cánh hỏng của mình, cố gắng bay lên trời với tư thế hết sức vụng về - cánh quăn queo, thân hình chao đảo - nhưng vẫn còn hơn là ngồi lên thứ đó. Trông cô như một con chim bị thương đang cố gắng bay, và điều đó khiến Marine phải lắc đầu.

\img{s-16e3}

Nàng Thuyền trưởng nhìn theo, thở dài rồi tặc lưỡi, vừa đạp xe vừa nói: ``Không leo thì thôi\dots\ Chứ chiếc xe này nhanh lắm nha! Chị sẽ chứng minh cho em thấy!''

Nói xong, cô đạp mạnh bàn đạp và phóng đi với tốc độ ``bàn thờ'' nhanh hơn cả cô nghĩ. Gió thổi tung mái tóc đỏ của cô, và cô có thể cảm nhận được sức mạnh của chiếc xe đang đưa cô đi xa hơn mỗi giây.

Cô phi thẳng về công ty với vẻ mặt đắc thắng, như thể vừa chứng minh rằng chiếc xe này không phải thứ bỏ đi. Cô rẽ trái, rẽ phải, len lỏi qua dòng người trên đường mà chẳng thèm giảm tốc độ, như một tay đua chuyên nghiệp trong một cuộc đua xe đạp địa hình.

``Chiếc xe này đúng là bá đạo! Chẳng ai đuổi kịp chị đâu!!''

Thế nhưng, khi đang chạy với vận tốc ngày càng tăng, Marine chợt nhận ra\dots\ phía trước là một khúc cua gắt với một dải phân cách cứng ngắc. Nó hiện ra như một bức tường bê tông, và cô không có đủ thời gian để phanh.

``Ể?''

Cô cố gắng hãm lại, nhưng, than ôi, cái xe này không có phanh để kìm lại được! Chân cô đạp vào không khí trong vô vọng, và chiếc xe cứ lao về phía trước như một mũi tên đã rời cung.

``\textbf{Ể!???}''

*RẦM!*

\img{s-16e4}

Một tiếng động khủng khiếp vang lên khi chiếc xe ba bánh đâm sầm vào dải phân cách, làm cho nàng Thuyền trưởng bật lên cao, xoay vài vòng trên không rồi đáp xuống đất như một bao cát. Bụi bay mù mịt, những người đi đường dừng lại nhìn, và những chiếc xe khác phải phanh gấp để tránh va chạm.

Marine thì\dots

``Khụ khụ\dots''

\dots cô nàng không sao cả, may mắn thay, không có lấy một vết thương. Dù đâm rất mạnh, cô vẫn đứng dậy phủi bụi, mặt không một vết trầy xước - cơ thể thuyền trưởng của cô dường như được làm bằng chất liệu cao su. Nhưng khi cô nhìn xuống chiếc xe của mình, cả cơ thể cô bỗng đơ lại.

Chiếc xe ba bánh hoàn toàn nát bươm.

Bánh trước văng ra một nơi xa tít, lăn lăn trên đường rồi biến mất sau một góc phố. Tay lái thì rụng xuống đất, bánh sau thì méo mó đến mức chẳng thể nhận ra hình dạng ban đầu. Những mảnh vỡ nằm rải rác khắp nơi, như một xác xe sau một vụ tai nạn thảm khốc.

``\textbf{Khôngggggggggggggggggggggggggg!!!!!!!!!!}''

Marine quỳ xuống đường, hai tay ôm đầu. Cô khóc lóc, dùng nắm đấm đập mạnh xuống đất, như thể đây là bi kịch lớn nhất đời cô, lớn hơn cả việc mất một kho báu hay bị Noel đá như một quả bóng.

\img{s-16e5}

Nhưng có một điều kỳ lạ.

Dù miệng cô há rộng, nước mắt lăn dài, và động tác trông có vẻ bi thảm, nhưng\dots

Không có tiếng khóc nào phát ra cả.

Marine đang khóc\dots\ không thành tiếng, theo đúng nghĩa đen. Mồm cô mở ra đóng vào như một con cá, nhưng không một âm thanh nào thoát ra, như thể giọng nói của cô đã bị nuốt chửng bởi nỗi đau.

Kanata, người đang bay lơ lửng trên trời với đôi cánh hỏng, nhìn xuống cảnh tượng này với vẻ mặt khó hiểu. Cô đưa tay che miệng, tự hỏi liệu mình có nên xuống giúp hay không.

Cùng lúc đó, Kuon cũng chứng kiến toàn bộ. Tôi nhìn Marine khóc lóc trong im lặng, rồi ngẩng lên nhìn Kanata đang bay loạng choạng trên trời.

Cả hai nhìn nhau một lúc, rồi đồng thanh thở dài.

Nàng Thiên sứ hỏi, giọng vừa thương cảm vừa buồn cười: ``\dots  Chị ấy đang làm cái quái gì thế? Em đã từng thấy nhiều người khóc, nhưng chưa thấy ai khóc mà không có tiếng cả.''

Cậu Tổng nhún vai, giọng đầy bất lực của một người đã chứng kiến quá nhiều chuyện: ``\dots  Anh cũng chịu, không hiểu nổi. Có lẽ đây là cách em ấy thể hiện nỗi đau trong im lặng.''

Marine vẫn quỳ đó, khóc không thành tiếng, đập đất như một đứa trẻ mất đồ chơi. Người qua đường nhìn cô với ánh mắt đầy thương cảm, nhưng cũng có người không nhịn được cười.

Về sau, chiếc xe đó được chính cậu con trai sửa lại - sau khi cậu thu thập tất cả các mảnh vỡ và nhờ một người bạn thợ sửa xe giúp - và được chính cô nàng Thuyền trưởng cảm kích đến mức cô đã tặng tôi một hộp sô-cô-la để cảm ơn (không có độc, cậu đã kiểm tra).

Nhưng, hình ảnh Marine quỳ trên đường, khóc không thành tiếng trước đống đổ nát của chiếc xe ba bánh lại là thứ kỳ quặc nhất mà cậu vẫn còn nhớ tới tận bây giờ. Có lẽ đó là lần đầu tiên cậu thấy cô ấy thực sự đau khổ vì một thứ không phải là kho báu hay tiền bạc.

\pagebreak

\subsection*{\phantomsection\label{ep-S16.VI}}
\addcontentsline{toc}{subsection}{{\sen{-Histoire 6-}} {\caviar Gập bụng}}
\stpt{-Histoire 6-}
\stcap{Gập bụng}
\vspace{-4pt}
{\caviar\fontsize{12pt}{12pt}\selectfont Nhân vật: \emp{kuon1} \emp{marine}}

Một buổi chiều, tại văn phòng, không khí yên tĩnh đến mức tôi có thể nghe thấy tiếng mình gõ bàn phím. Bên kia bàn, Marine ngồi lướt điện thoại đọc bình luận trên kênh của mình, thỉnh thoảng cười khúc khích vì những bình luận hài hước. Nhưng rồi, nụ cười của cô tắt dần. Mắt cô giật giật, và tôi có thể cảm nhận được nhiệt độ trong phòng đang tăng lên.

Bỗng nhiên, cô đập tay xuống bàn một cái *RẦM*, khiến tôi giật mình suýt làm đổ cốc nước: ``Kuon-senpai! Quay video giúp em!''

Tôi ngơ ngác, tay vẫn còn trên bàn phím: ``\dots  Hả? Là sao? Quay\dots\ gì cơ?''

Em ấy không trả lời, cô đứng dậy, bước ra giữa phòng với vẻ mặt đầy quyết tâm của một chiến binh chuẩn bị bước vào trận chiến. Tôi vẫn chưa biết chuyện gì đang xảy ra, nhưng tôi đã học được rằng ở công ty này, tốt nhất là không nên hỏi quá nhiều.

\ct

Tôi kiểm tra lại camera trên điện thoại, rồi nhìn lên, thấy Marine đã đứng sẵn ở giữa phòng với vẻ mặt vô cùng nghiêm túc còn hơn cả lúc cô tính toán chia kho báu: ``Chuẩn bị máy xong chưa, Kuon-senpai!?''

Tôi giơ ngón tay cái lên: ``Bấm máy rồi đây. Sẵn sàng khi em sẵn sàng.''

Tôi vẫn chưa hiểu tại sao cô nàng này lại đột nhiên muốn quay một thước phim tập luyện thể thao. Nhưng dù sao cũng không có lý do gì để từ chối - biết đâu lại có nội dung hay cho kênh của cô nàng thì sao.

``OK\dots\ sẵn sàng\dots\ Bắt đầu!''

\il

Marine giơ tay chào màn hình với nụ cười rạng rỡ - cái nụ cười mà cô thường dùng để chào đón khán giả trong các buổi stream: ``Ahoy! Thuyền trưởng Marine đây! Hôm nay ta sẽ tập một bài gập bụng để rèn luyện sức khỏe và sức lực nha!''

Tôi nhướng mày, cười khẩy trong đầu. Cô ấy có vẻ nghiêm túc, nhưng tôi đã thấy nhiều người hùng hục tập thể dục rồi bỏ cuộc sau năm phút. Nàng Thuyền trưởng hùng hồn tuyên bố: ``Chúng ta cùng bắt đầu nào!'' đồng thời giơ nắm đấm lên cao, như thể đang chuẩn bị cho một trận đấu boxing.

\ct

Cảnh chuyển đến Marine đã vào tư thế chuẩn bị. Cô đã cởi mũ, bịt mắt và áo choàng ra, chỉ còn lại bộ đồ hóa trang làm hải tặc quen thuộc - trông cô lúc này như một vận động viên thể dục dụng cụ đang chuẩn bị bước vào sân khấu Olympic, chỉ khác là sân khấu của cô là sàn nhà văn phòng.

Tôi đứng sau máy quay, khoanh tay, giọng đầy hoài nghi: ``Chả biết có ra ngô ra khoai gì không đây\dots\ Nhưng thôi, cứ quay đã.''

Marine hít một hơi sâu, dồn toàn bộ sức lực vào cơ bụng. Cô nằm xuống sàn, hai tay đặt sau đầu, rồi bắt đầu cố gắng nâng người lên. Mặt cô đỏ bừng, mồ hôi bắt đầu lấm tấm trên trán.

Tôi đếm, giọng đều đều như một máy đếm tự động: ``\dots  Một.''

Nhưng tôi nhíu mày khi thấy cô chỉ nâng được nửa người, lưng vẫn chưa chạm tới đầu gối: ``Vẫn chưa ngồi dậy hẳn đâu. Đấy mới chỉ là nửa cái gập bụng.''

\img{s-16f1}

Marine nghiến răng, dồn thêm sức để gập lần thứ hai. Nhưng lần này, cô gần như không thể nâng nổi. Cơ thể cô run lên bần bật, và một tiếng rên vang lên từ sâu trong cổ họng: ``Aaa\dots\ aaaaaa\dots\ \textbf{AAAAHH\dots!!}''

Tiếng rên của cô vang khắp phòng, kéo dài, tha thiết, và nghe giống như\dots\ đang đỡ đẻ hơn là tập thể dục. Tôi chớp mắt, toát mồ hôi hột, tự hỏi liệu mình có đang quay một video thể thao hay một bộ phim tài liệu về sinh nở: ``\dots  Hai.''

Ngay sau đó, Marine rơi bịch xuống sàn như một con cá chết, thở phì phò, mồ hôi vã ra như tắm. Toàn thân cô nằm bẹp trên sàn, hai tay duỗi thẳng, mắt nhìn lên trần nhà với vẻ mặt như vừa hoàn thành một cuộc chạy marathon trăm cây số.

Tôi cúi xuống nhìn, giọng đầy ngạc nhiên: ``\dots  Mới được có hai cái mà đã thở phì phò rồi. Sao hôm nay yếu thế? Em thường tự hào về sức bền của mình mà.''

Tôi nheo mắt, rồi chợt nghĩ ra một giả thuyết, một giả thuyết mà tôi không thể giữ trong đầu: ``\dots  Hay là\dots\ định thở một cách dâm đãng để hút người xem? Đây là chiến lược nội dung mới của em à?''

\img{s-16f2}

Marine nằm bẹp trên sàn, không còn đủ sức để cãi lại. Cô chỉ có thể đưa mắt về phía máy quay, cố tỏ ra chuyên nghiệp như thể đây là kết thúc của một video đầy thành công. Mặc dù cô thở hổn hển, vẫn cố gắng mỉm cười và giơ ngón tay cái lên, như một vận động viên vừa hoàn thành bài thi dù biết rằng mình đã thất bại thảm hại.

Tôi tắt máy quay, nhìn cô nằm bất động trên sàn, và thở dài: ``Anh nghĩ là em nên bắt đầu bằng những bài tập nhẹ hơn. Gập bụng hai cái trong mười phút - có lẽ nên tập hít thở trước đã.''

Hình ảnh dần trắng xóa, và đoạn phim kết thúc tại đó.

Về sau, Marine không bao giờ nhắc lại video này, và tôi không bao giờ hỏi. Nhưng mỗi khi thấy cô ấy cầm điện thoại lên, tôi lại tự hỏi liệu cô ấy có đang đọc bình luận về ``già'' hay ``xuống cấp'' nữa không.

\subsection*{\phantomsection\label{ep-S16.VII}}
\addcontentsline{toc}{subsection}{{\sen{-Histoire 7-}} {\caviar Chống đẩy}}
\stpt{-Histoire 7-}
\stcap{Chống đẩy}
\vspace{-4pt}
{\caviar\fontsize{12pt}{12pt}\selectfont Nhân vật: \emp{kuon1} \emp{marine}}

Một ngày sau, video gập bụng của Marine đã được đăng tải lên mạng.

Và tất nhiên, phản ứng của khán giả không như cô mong đợi. Thay vì những lời cổ vũ, cô nhận được vô số những lời trêu chọc và châm biếm, thậm chí còn thê thảm hơn, dạng như:

``Ủa? Chị gập bụng có hai cái thôi hả?''

``Thở như vậy có ổn không? Hay là tập xong chị cần đi khám bệnh luôn?''

``Bà già Marine yếu quá đi!''

Nhìn những bình luận châm chọc từ mọi người, tôi cười trừ, như đã xác định được trước: ``Anh không ngạc nhiên. Mới có gập được hai cái thôi mà đã thở rồi đấy, bảo sao\dots\ Em muốn chứng minh mình không già thì ít nhất cũng phải làm được mười cái.''

Marine nghiến răng. Cô đập tay xuống bàn, đứng bật dậy như một lò xo bị nén quá mức: ``Không thể để như thế này được!''

Thế là hôm nay, em lại nhờ tôi quay tiếp một video tập thể thao khác. Lần này em quyết tâm chứng minh rằng mình không hề yếu đuối - em sẽ làm được ít nhất hai mươi cái chống đẩy.

Nhưng lần này, tôi quyết định đẩy trách nhiệm quay phim cho Zentarou - vì hôm nay cậu ta đang rảnh rỗi và tôi thì không muốn mất thêm thời gian cho những trò điên rồ của Marine. Còn bản thân tôi, tôi sẽ chỉ đứng giám sát.

\ct

Marine đứng ở giữa phòng, hai tay chống nạnh, vẻ mặt quyết tâm. Zentarou đã sẵn sàng với máy quay, tay cầm chắc như một nhiếp ảnh gia chuyên nghiệp: ``Máy quay sẵn sàng rồi chứ?''

``OK-la rồi! Chuẩn bị\dots''

Tôi khoanh tay, gật đầu, giọng đều đều: ``\textbf{Bắt đầu!}''

\il

Marine tươi cười nhìn vào camera, vẫn với bộ trang phục như lần trước, nhưng lần này cô có thêm một chiếc băng đô màu đỏ trên trán, trông như một vận động viên thể dục thập niên 80: ``Ahoy! Thuyền trưởng Marine đây! Vì mình không muốn sau này bị viêm khớp nên hôm nay mình sẽ tập thể lực nha! Mọi người hãy cùng cổ vũ cho mình!''

Tôi đứng ngoài khung hình, lẩm bẩm đủ để Zentarou nghe thấy: ``Để xem một bà già `mãi mãi tuổi 17' này làm gì nào. Tôi cá là cô ấy sẽ không qua nổi mười cái.''

Zentarou nghiêng đầu, giọng đầy tò mò: ``\vie{Ý mày là sao, Kuon? Cô ấy trông có vẻ nghiêm túc mà.}''

``\vie{Cứ xem đi rồi sẽ thấy. Tao đã chứng kiến quá nhiều lần rồi.}''

Marine giơ tay đầy khí thế, như một vị tướng ra lệnh tấn công: ``Giờ thì, ra khơi nào!''

\ct

Giống như lần trước, cô đã cởi mũ, bịt mắt và áo choàng ra để tiện tập luyện. Cô đã ở tư thế chuẩn bị chống đẩy - hai tay đặt trên sàn, lưng thẳng, cơ bụng căng lên.

``Sẵn sàng chưa? Bắt đầu nào!''

\img{s-16g1}

Zentarou chỉnh lại góc quay, hướng về phía Marine đang chống đẩy. Cậu ta có vẻ tập trung, nhưng\dots

Nhưng ngay khi Marine bắt đầu hạ người xuống\dots

Ống kính máy quay từ từ\dots\ hạ xuống theo.

Cụ thể hơn, nó tập trung vào phần ngực của cô. Từ góc này, từ góc kia, Zentarou quay như một đạo diễn đang thực hiện một bộ phim tài liệu về ``vòng một của thuyền trưởng''.

Tôi lập tức nhíu mày, giọng đầy bất mãn: ``\vie{Zentarou-senpai, mày làm cái vẹo gì đấy? Quay cho tử tế vào chứ. Đây là video thể thao, không phải phim 18+!}''

Nhưng cậu ta không hề đáp lại. Thay vào đó, cậu ta tiếp tục lia máy quay ở nhiều góc độ khác nhau, chủ yếu tập trung vào vòng một của Thuyền trưởng, như thể đang ghi hình cho một buổi chụp ảnh thời trang nội y.

Tôi ngán ngẩm, miệng vẫn đếm số cái chống đẩy của Marine, nhưng mắt thì lia về khuôn mặt đầy vẻ dâm đãng của cậu bạn: ``Một. Hai\dots\ \vie{Zentarou, mày đùa với tao đấy à? Tao bảo mày quay thể thao, không phải quay kiểu này.}''

Zentarou vẫn im lặng, chỉ lẩm bẩm với vẻ mặt đầy đạo đức giả, như thể đang thực hiện một công việc khoa học quan trọng: ``\vie{Quay góc này nhìn đẹp ghê. Nó rất\dots\ nghệ thuật.}''

Tôi nhìn thấy cảnh này, chỉ biết day trán. Mồ hôi lấm tấm trên trán tôi, và tôi thầm hối hận vì đã để cậu ta cầm máy quay: ``\vie{Mày dâm đãng đến thế luôn\dots\ Mày đang làm xấu mặt tao đấy\dots}''

\ct

Tôi tiếp tục đếm, cố gắng tập trung vào Marine để không nhìn vào màn hình máy quay: ``Bảy\dots\ Tám\dots''

Đến đây, Marine bắt đầu lung lay, tay run lên bần bật. Cơ thể cô như một chiếc thuyền trong cơn bão, và cuối cùng, cô nằm bẹp xuống sàn, thở hổn hển, mồ hôi vã ra như tắm.

Zentarou đặt máy quay xuống, hướng về khuôn mặt đầy mệt mỏi của cô, vẫn giữ một nụ cười mờ ám trên môi.

\img{s-16g2}

Tôi khoanh tay, nhận xét với giọng của một huấn luyện viên khó tính: ``Tám cái chống đẩy. Kém. Còn tệ hơn cả lần gập bụng trước.''

Zentarou gật gù, vẫn không rời mắt khỏi màn hình máy quay: ``Công nhận. Nhưng mà tôi vẫn quay được mấy góc đẹp.''

Marine thở một lúc lâu, rồi đột ngột ngẩng lên, dí sát mặt vào ống kính. Mắt cô trợn tròn, giọng như một cơn bão sắp ập đến: ``Này mấy ông kia! Đúng, mấy ông đang nhìn vào ngực của tôi đấy! Tôi biết hết!''

\img{s-16g3}

Tôi nhếch môi, liếc qua Zentarou với ánh mắt đầy bất lực: ``Gòi xong. Tại mày đấy, senpai ơi. Giờ thì mày chịu trận đi.''

Zentarou ho nhẹ, giả vờ không liên quan, mắt nhìn lên trần nhà như thể đang ngắm mây trời: ``Ủa a-lô? Tao không nghe thấy gì cả.''

Marine quay sang Zentarou, chống nạnh, giọng như một vị thuyền trưởng ra lệnh trừng phạt thủy thủ đoàn: ``Chống đẩy 100 cái cho tôi nhé!''

Cậu bạn Bách khoa giật mình, mắt mở to. Cô chỉ thẳng vào cậu ta, như thể đang chỉ định một tội phạm: ``Anh đó, còn nhìn ai nữa! Cứ coi như đó là hình phạt cho việc quay lung tung của anh đi!''

Zentarou chớp mắt, rồi ngay lập tức bật chế độ chối bay chối biến, giơ hai tay lên như một kẻ đầu hàng: ``Ơ, anh không biết, không biết gì nha! Anh chỉ quay theo yêu cầu thôi!''

Rồi cậu ta\dots\ vội vã chạy khỏi phòng, để lại Marine đứng đó với vẻ mặt đắc thắng, và tôi đứng đó với vẻ mặt bất lực.

Tôi nhìn theo bóng dáng Zentarou biến mất, rồi quay sang Marine, thở dài: ``Anh nghĩ em nên tập trung vào việc cải thiện số cái chống đẩy thay vì đuổi người quay phim đi.''

Marine cười khúc khích, vỗ vai tôi: ``Thôi kệ, cứ để cậu ta chạy đi. Lần sau em sẽ tự tập một mình.''

Tôi lắc đầu, vừa cười vừa thở dài: ``Anh chỉ hy vọng là lần sau em sẽ tập được nhiều hơn tám cái.''

Đoạn phim kết thúc tại đó, với hình ảnh Marine đứng giữa phòng, tay chống nạnh, vẻ mặt đắc thắng của một người vừa chiến thắng một trận chiến nhỏ - dù cô đã thua trong trận chiến chống đẩy.

Tôi tắt máy quay, và thầm nghĩ: ``\textit{Ít nhất thì hôm nay cũng có một người bị bắt phạt. Dù không phải là người đáng bị phạt.}''

\pagebreak

\subsection*{\phantomsection\label{ep-S16.VIII}}
\addcontentsline{toc}{subsection}{{\sen{-Histoire 8-}} {\caviar Chạy vượt rào}}
\stpt{-Histoire 8-}
\stcap{Chạy vượt rào}
\vspace{-4pt}
{\caviar\fontsize{12pt}{12pt}\selectfont Nhân vật: \emp{kuon1} \emp{marine}}

Ba ngày, ba thước phim thể thao liên tiếp.

Và mỗi lần như vậy, số lượng bình luận trêu chọc Marine chỉ có tăng chứ không giảm. Tôi đã đọc qua vài bình luận, và tôi có thể thấy rằng khán giả của cô ấy không hề nương tay - họ châm chọc cô như một môn thể thao mới.

``Có nhất thiết phải hơn thua mọi người như thế không, Marine?'' tôi ngán ngẩm hỏi, chuyển bị máy quay phim từ tay này sang tay kia, tự hỏi mình đang làm cái quái gì với cuộc đời này.

Tokue cười trừ, giọng đầy thương hại nhưng cũng không kém phần châm biếm: ``Marine à, tôi nghĩ bà nên từ bỏ đi. Lần sau quay video tập dưỡng sinh đi, đừng cố quá nữa, kẻo thành quá cố đấy. Bà tập thể thao không phải sở trường của bà.''

Marine nghiến răng với những lời châm chọc không chỉ từ những người xem của cô, mà còn chính từ đồng nghiệp của mình. Tôi có thể thấy đôi mắt cô ánh lên tia lửa - không phải lửa giận, mà là lửa của sự kiên trì.

Không thể để như vậy được. Cô quyết định làm một video thể thao nữa - và lần này, là chạy vượt rào.

Tôi chỉ biết thở dài, đưa tay xoa trán: ``Mình phải làm như thế này bao nhiêu lần nữa đây\dots? Sáng giờ chưa uống được ngụm cà phê nào.''

\ct

Vì không muốn lặp lại cảnh tượng hôm trước - khi Zentarou đã quay một video tập trung vào ngực Marine thay vì tập thể dục - lần này tôi không nhờ cậu ta nữa. Tôi kéo theo Tokue, người ít ra không ``dâm dê'' bằng cậu bạn Bách khoa kia. Ít nhất thì tôi tin rằng cậu ta sẽ quay đúng hướng.

Lúc cả ba đến nơi, Tokue tò mò hỏi, giọng đầy vẻ không hiểu nổi: ``Biết rõ là `cụ bà' Marine không khỏe rồi, vậy mà sao lại cố làm video thể thao thế? Chẳng lẽ bà ấy không biết mình đang tự hành hạ mình?''

Tôi chỉ tặc lưỡi, giọng đầy vẻ bất lực của một người đã từ bỏ việc giải thích: ``Kệ em ấy đi. Thích thể hiện thì cứ làm cho em ấy thôi. Dù sao thì cũng là nội dung cho kênh của em ấy.''

Họ chọn một cung đường rộng, vắng vẻ để quay - một con đường dài khoảng trăm mét, không có xe cộ qua lại. Marine đã tự chuẩn bị ba rào chắn để làm chướng ngại vật cho phần thi của mình. Trông chúng như những rào chắn thật, nhưng tôi nghi ngờ chúng được làm từ ống nhựa và băng keo.

Nhìn vào số rào chắn, tôi và Tokue chỉ lẳng lặng nhìn nhau.

Không cần nói ra, cả hai đều có chung suy nghĩ: ``\textit{Kiểu gì cũng ngã ở rào đầu tiên cho mà xem.}''

\il

Khác với hai lần trước, lần này Marine không nói gì cả. Cô không giơ tay chào, không hô ``Ahoy!'', không nói một lời nào. Cô chỉ đứng sẵn ở tư thế chuẩn bị, mắt nhìn thẳng về phía trước, như một vận động viên Olympic sắp bước vào cuộc thi.

Máy quay (do tôi cầm) hướng xuống cô. Có thể thấy, khuôn mặt Marine đang thể hiện một vẻ tràn đầy quyết tâm - kiểu quyết tâm mà tôi chỉ thấy khi cô ấy đang tính toán chia kho báu.

Trước mặt cô là ba rào chắn, cách nhau đều đặn trong khoảng năm mươi mét. Mỗi rào cao khoảng nửa mét, đủ để một người bình thường phải nhảy qua, nhưng với Marine thì\dots

\img{s-16h1}

Tokue cầm đồng hồ bấm giờ, giọng như một trọng tài chuyên nghiệp: ``Chuẩn bị\dots''

``Sẵn sàng\dots''

``\textbf{Chạy!}''

Vừa dứt lời, Marine lao về phía trước. Tốc độ của cô nhanh hơn tôi tưởng - có lẽ vì cô đã tập trung tất cả sức mạnh vào đôi chân.

Ngay lập tức, cô đối mặt với chướng ngại đầu tiên.

Tôi và Tokue nheo mắt, sẵn sàng chứng kiến cảnh Marine cố nhảy qua rào - và tất nhiên, chúng tôi dự đoán cô sẽ ngã.

Nhưng không.

Thay vì nhảy\dots

Marine xô đổ cả cái rào xuống, và bất ngờ thay, cô nàng lại tiếp tục băng băng về phía trước! Cô không hề dừng lại, không hề ngần ngại, như thể đó là cách vượt rào đúng đắn.

\img{s-16h2}

Tokue trợn mắt, giọng như sắp vỡ: ``Oắt-đờ-''

Tôi sững sờ, tay cầm máy quay run lên: ``Cái quái\dots? Đó có phải là vượt rào không hả?''

Cái thứ hai? Xô đổ.

Cái thứ ba? Lật úp luôn.

Dù không vấp ngã lần nào, nhưng cách cô ``vượt rào'' này - xô đổ tất cả các chướng ngại vật thay vì nhảy qua - làm hai chúng tôi chỉ biết đứng hình, không biết nên phản ứng như thế nào.

Cuối cùng, sau khi chạy qua ba chướng ngại, Marine dừng lại trước máy quay, thở dốc, mồ hôi nhễ nhại.

Tôi khoanh tay, nhướng mày, giọng đầy ngạc nhiên: ``Làm gì mà thở phì phò như trâu thế? Mới tí mà đã mệt rồi? Em chỉ mới chạy có năm mươi mét thôi.''

Nhưng, nhân vật chính của thước phim không đáp lại. Cô chỉ ngẩng lên nhìn vào máy quay. Mồ hôi nhễ nhại, mái tóc đỏ bết vào trán. Nhưng vẫn giơ ngón tay cái lên, như thể muốn nói: ``Tôi làm được rồi nhé mọi người\~{}''

\img{s-16h3}

Cả tôi và Tokue đồng thanh, giọng như một bản hòa ca bất đắc dĩ: ``Có thành công đíu đâu mà tỏ ra tự hào thế\dots\ Em vừa xô đổ tất cả rào chắn, chứ có nhảy qua cái nào đâu.'' / ``Đó có phải là chạy vượt rào đâu mẹ trẻ! Đó là chạy đập rào!''

\il

Ba video đã được đăng lên mạng.

Và phản ứng của khán giả? Như cả tôi, Zentarou và Tokue dự đoán, khen thì ít, mà chê (hay châm chọc) thì nhiều.

Tôi lướt qua các bình luận và thấy những dòng như:

``Marine à\dots\ tôi không biết nên khen hay chê bà nữa.''

``Cứng đầu vcl''

``Có ai đó ngăn bà già này lại được không?''

``Thôi bỏ đi chị ơi, chạy bộ bình thường thôi cũng là cả một thành tích rồi!''

``Xô rào chứ có nhảy đâu mà gọi là vượt rào hả trời?''

Những bình luận kiểu như vậy khiến cho Marine mất hoàn toàn động lực. Cô ngồi xó góc, đầu rụp xuống, ôm mặt khóc rưng rức. Lần này là khóc thật, không phải khóc không tiếng như vụ xe ba bánh.

Trong khi đó, tôi ngồi trên ghế, vươn vai đầy mệt mỏi, mắt nhìn về phía cô nàng đang đau khổ: ``\vie{Có khi nào mai em ấy lại đòi quay nữa không? Tôi không chịu nổi nữa đâu.}''

Tokue lắc đầu, giọng đầy thương cảm: ``\vie{Tôi không nghĩ vậy đâu. Nhìn bà ấy mất hết tinh thần rồi kia kìa. Rõ từ bỏ là cái chắc.}''

Zentarou thì cười hề hề, vẫn chưa từ bỏ ý định được cầm máy quay: ``\vie{Nói chứ, nếu mai quay nữa, cho tao đi quay tiếp nha. Tao hứa sẽ quay\dots\ nghệ thuật.}''

Tôi nhìn cậu bạn Bách khoa của mình với ánh mắt đầy sát khí, giọng như một lời cảnh cáo cuối cùng: ``\vie{Éo. Riêng mày thì cấm động vào máy quay luôn. Tao không muốn thêm một vụ kiện nào nữa.}''

Marine vẫn ngồi đó, ôm mặt, nước mắt lăn dài. Nhưng tôi có thể thấy đôi mắt cô vẫn ánh lên một tia gì đó - không phải sự bỏ cuộc, mà là sự kiên trì.

Tôi nhìn cô, rồi thầm nghĩ: ``\textit{Chắc chắn là mai em ấy sẽ lại đòi quay tiếp. Nhưng lần sau thì cứ nói là máy quay hỏng cho khỏi quay nữa, chứ cứ thế này càng quay càng tự xúc đất đào đất chôn mình xuống rồi còn gì\dots}''

\end{document}